\documentclass[../main.tex]{subfiles}
\begin{document}
\chapter{拓扑空间}\label{chapter:appA-拓扑空间}

为读者方便，此处收集若干拓扑学概念与结果（例如见文献~\cite{ref60}）。

\section{拓扑与距离空间}\label{section:appA-拓扑定义}

\begin{definition}[拓扑与拓扑空间]\label{definition:appA-拓扑}
设 $X$ 为集合，$\mathscr{T}\subseteq\mathcal{P}(X)$。称 $\mathscr{T}$ 为 $X$ 上的一个拓扑，若：
\begin{enumerate}[(i)]
\item $\emptyset\in\mathscr{T}$，$X\in\mathscr{T}$；
\item 对任意族 $\{A_i\}_{i\in I}\subseteq\mathscr{T}$，有 $\bigcup_{i\in I}A_i\in\mathscr{T}$；
\item 对任意有限族 $\{A_1,\dots,A_n\}\subseteq\mathscr{T}$，有 $\bigcap_{k=1}^n A_k\in\mathscr{T}$。
\end{enumerate}
偶 $(X,\mathscr{T})$ 称为拓扑空间。集合 $A\in\mathscr{T}$ 称为开集；若 $C^c\in\mathscr{T}$，则称 $C\subseteq X$ 为闭集。
\end{definition}

\begin{definition}[强弱拓扑、生成拓扑与相对拓扑]\label{definition:appA-拓扑关系}
若 $\mathscr{T}_1,\mathscr{T}_2$ 均为 $X$ 上的拓扑且 $\mathscr{T}_1\subseteq\mathscr{T}_2$，则称 $\mathscr{T}_1$ 弱于 $\mathscr{T}_2$，或称 $\mathscr{T}_2$ 强于 $\mathscr{T}_1$。拓扑 $\mathcal{P}(X)$ 称为离散拓扑。对任意族 $\mathscr{G}\subseteq\mathcal{P}(X)$，交
\begin{align*}
\mathscr{T}_0(\mathscr{G}):=\bigcap\{\mathscr{T}\text{ 为 }X\text{ 上的拓扑}\mid\mathscr{G}\subseteq\mathscr{T}\}
\end{align*}
是一个拓扑，称为由 $\mathscr{G}$ 生成的拓扑，或称包含 $\mathscr{G}$ 的最小拓扑。对任意 $F\subseteq X$，族
\begin{align*}
\mathscr{T}|_F:=\{A\cap F\mid A\in\mathscr{T}\}
\end{align*}
是 $F$ 上的拓扑，称为 $\mathscr{T}$ 在 $F$ 上的相对拓扑或迹；拓扑空间 $(F,\mathscr{T}|_F)$ 称为 $(X,\mathscr{T})$ 的拓扑子空间。
\end{definition}

\begin{definition}[距离拓扑]\label{definition:appA-距离拓扑}
设 $(X,d)$ 为距离空间，$d$ 为其上的距离。考虑开球族
\begin{align*}
B(x_0,r):=\{x\in X\mid d(x_0,x)<r\}\qquad(x_0\in X,\ r>0),
\end{align*}
集合族
\begin{align*}
\{E\subseteq X\mid \forall x_0\in E,\ \exists r>0\text{ 使得 }B(x_0,r)\subseteq E\}
\end{align*}
是 $X$ 上的一个拓扑，称为距离拓扑（因而每个距离空间都是拓扑空间）。
\end{definition}

\begin{definition}[可度量化]\label{definition:appA-可度量化}\label{proposition:appA-可度量化}
拓扑空间 $(X,\mathscr{T})$ 称为：
\begin{enumerate}[(i)]
\item 可度量化的，若存在 $X$ 上的距离 $d$，使相应距离拓扑与 $\mathscr{T}$ 一致；
\item 完全可度量化的，若存在 $X$ 上的距离 $d$，使相应距离拓扑与 $\mathscr{T}$ 一致，且 $(X,d)$ 为完备距离空间。
\end{enumerate}
\end{definition}

\begin{definition}[点集间距离、直径与有界集]\label{definition:appA-距离直径}
设 $(X,d)$ 为距离空间。对任意非空 $E,F\subseteq X$，$E$ 与 $F$ 之间的距离定义为
\begin{align*}
d(E,F):=\inf_{x\in E,\,y\in F}d(x,y),
\end{align*}
点 $x\in X$ 到非空集 $E\subseteq X$ 的距离为 $d(x,E):=d(\{x\},E)$。按约定，对所有 $x\in X$ 有 $d(x,\emptyset):=\infty$。非空集 $E\subseteq X$ 的直径为
\begin{align*}
\operatorname{diam}(E):=\sup_{x,y\in E}d(x,y)\in[0,\infty],
\end{align*}
并约定 $\operatorname{diam}(\emptyset):=-\infty$。若 $\operatorname{diam}(E)<\infty$，则称 $E\subseteq X$ 为有界集。
\end{definition}

\section{乘积拓扑}\label{section:appA-乘积拓扑}

\begin{definition}[乘积拓扑]\label{definition:appA-乘积拓扑}
设 $(X_1,\mathscr{T}_1)$ 与 $(X_2,\mathscr{T}_2)$ 为拓扑空间。由族
\begin{align*}
\{A_1\times A_2\mid A_1\in\mathscr{T}_1,\ A_2\in\mathscr{T}_2\}\subseteq\mathcal{P}(X_1\times X_2)
\end{align*}
生成的拓扑称为乘积拓扑，记作 $\mathscr{T}_1\times\mathscr{T}_2$。拓扑空间 $(X_1\times X_2,\mathscr{T}_1\times\mathscr{T}_2)$ 称为 $X_1$ 与 $X_2$ 的拓扑乘积。
\end{definition}

\section{邻域与 Hausdorff 空间}\label{section:appA-邻域Hausdorff}

\begin{definition}[邻域]\label{definition:appA-邻域}
设 $(X,\mathscr{T})$ 为拓扑空间。集合 $E\subseteq X$ 的邻域是指任意满足"存在 $A\in\mathscr{T}$ 使 $E\subseteq A\subseteq F$"的集合 $F\subseteq X$。单点集 $\{x\}$ 的邻域称为点 $x$ 的邻域。
\end{definition}

\begin{definition}[Hausdorff 空间]\label{definition:appA-Hausdorff}
拓扑空间 $X$ 称为 Hausdorff 空间，若对任意 $x_1,x_2\in X$，$x_1\neq x_2$，存在 $x_1$ 的邻域 $E_1$ 与 $x_2$ 的邻域 $E_2$，使 $E_1\cap E_2=\emptyset$。每个距离空间都是 Hausdorff 空间。
\end{definition}

\section{基、内部、闭包与可分性}\label{section:appA-基闭包可分}

\begin{definition}[拓扑的基]\label{definition:appA-基}
设 $(X,\mathscr{T})$ 为拓扑空间。族 $\{A_i\}_{i\in I}\subseteq\mathscr{T}$ 称为 $\mathscr{T}$ 的一个基，若对每个 $A\in\mathscr{T}$，存在 $I_0\subseteq I$ 使 $A=\bigcup_{i\in I_0}A_i$。若 $X$ 为距离空间，则开球族 $\{B(x_0,r)\mid x_0\in X,\ r>0\}$ 是距离拓扑的一个基。
\end{definition}

\begin{definition}[内部、闭包与边界]\label{definition:appA-内部闭包边界}
设 $E\subseteq X$。集合
\begin{align*}
E^\circ:=\bigcup\{A\in\mathscr{T}\mid A\subseteq E\}
\end{align*}
称为 $E$ 的内部。$E$ 的闭包定义为
\begin{align*}
\overline{E}:=\bigcap\{C\supseteq E\mid C^c\in\mathscr{T}\}.
\end{align*}
$E$ 的边界定义为 $\partial E:=\overline{E}\setminus E^\circ$。
\end{definition}

\begin{definition}[稠密与可分]\label{definition:appA-稠密可分}
设 $E,F\subseteq X$。若 $F\subseteq\overline{E}$，则称 $E$ 在 $F$ 中稠密；特别地，若 $\overline{E}=X$，则称 $E$ 在 $X$ 中稠密。若存在可数稠密集 $E\subseteq X$，则称空间 $X$ 为可分的。
\end{definition}

\begin{proposition}\label{proposition:appA-基可分}\label{proposition:appA-紧可分}
\begin{enumerate}[(i)]
\item 具有可数基的拓扑空间是可分的；
\item 可分距离空间具有可数基。
\end{enumerate}
\end{proposition}

\section{紧性}\label{section:appA-紧性}

\begin{definition}[开覆盖与紧集]\label{definition:appA-紧性}
设 $(X,\mathscr{T})$ 为拓扑空间。族 $\{A_i\}_{i\in I}\subseteq\mathscr{T}$ 称为集合 $E\subseteq X$ 的开覆盖，若 $E\subseteq\bigcup_{i\in I}A_i$。若 $X$ 的每个开覆盖都含有限子覆盖，则称 $X$ 为紧空间。子集 $K\subseteq X$ 称为紧集，若拓扑空间 $(K,\mathscr{T}|_K)$ 是紧的；若闭包 $\overline{K}$ 是紧的，则称 $K$ 为相对紧集。
\end{definition}

紧空间的每个闭子集是紧的；Hausdorff 空间 $X$ 的每个紧子集是闭的。若 $X_1$ 与 $X_2$ 为紧空间，则拓扑乘积 $X_1\times X_2$ 是紧的。

关于 $\mathbb{R}^n$ 中闭球对集合的覆盖，有如下定义。

\begin{definition}[细覆盖]\label{definition:appA-细覆盖}
闭球族 $\mathscr{G}:=\{B_r(x_0)\mid x_0\in\mathbb{R}^n,\ r>0\}$ 称为集合 $U\subseteq\mathbb{R}^n$ 的细覆盖，若对所有 $x_0\in U$ 有
\begin{align*}
\inf\{r\mid B_r(x_0)\in\mathscr{G}\}=0.
\end{align*}
\end{definition}

\begin{definition}[收敛、聚点与完备距离空间]\label{definition:appA-收敛聚点}
设 $(X,\mathscr{T})$ 为拓扑空间。序列 $\{x_n\}\subseteq X$ 收敛于 $x_0\in X$，若对 $x_0$ 的每个邻域 $F$，存在 $\bar n\in\mathbb{N}$，使当 $n>\bar n$ 时 $x_n\in F$；若 $x_0$ 的每个邻域都含有 $\{x_n\}$ 中无穷多个元素，则称 $x_0$ 为 $\{x_n\}$ 的聚点。若距离空间中每个 Cauchy 序列 $\{x_n\}\subseteq X$ 都收敛于某个 $x\in X$，则称该距离空间是完备的。
\end{definition}

\begin{definition}[可数紧与序列紧]\label{definition:appA-可数序列紧}
拓扑空间 $(X,\mathscr{T})$ 称为可数紧的，若 $X$ 的每个可数开覆盖都含有限子覆盖；$X$ 称为序列紧的，若每个序列 $\{x_n\}\subseteq X$ 都含收敛子列。子集 $F\subseteq X$ 称为序列紧的，若拓扑空间 $(F,\mathscr{T}|_F)$ 是序列紧的；若闭包 $\overline{F}$ 是序列紧的，则称 $F$ 为相对序列紧的。空间 $X$ 为可数紧当且仅当每个 $\{x_n\}\subseteq X$ 都有聚点。
\end{definition}

若 $X$ 为距离空间，则下述陈述等价：(i) $X$ 紧；(ii) $X$ 可数紧；(iii) $X$ 序列紧。距离空间的每个相对序列紧子集是有界的。

\section{局部紧与 Polish 空间}\label{section:appA-局部紧Polish}\label{subsection:appA-Polish}

\begin{definition}[局部紧与 $\sigma$-紧]\label{definition:appA-局部紧}
拓扑空间 $X$ 称为局部紧的，若每个 $x\in X$ 都有一个紧邻域。局部紧 Hausdorff 空间若为可数个紧集之并，则称为 $\sigma$-紧的。
\end{definition}

\begin{remark}\label{remark:appA-可数基局部紧}
每个具有可数基的局部紧 Hausdorff 空间 $X$ 的所有开子集都是 $\sigma$-紧的。事实上，设 $\{A_n\}\subseteq\mathscr{T}$ 为可数基，$A\in\mathscr{T}$。对每个 $x\in A$，存在 $x$ 的紧邻域 $K_x\subset A$，因而对某个 $k\in\mathbb{N}$ 有
\begin{align*}
x\in\bigcup_{n=1}^k A_n\subset\overline{\bigcup_{n=1}^k A_n}\subset K_x.
\end{align*}
显然每个集合 $\overline{\bigcup_{n=1}^k A_n}$ 是紧的；此外这些集合构成的族是可数的，且 $A$ 等于它们的并。
\end{remark}

与\cref{remark:appA-可数基局部紧}相一致，可证明如下结果。

\begin{proposition}\label{proposition:appA-局部紧基}\label{proposition:appA-σ紧}\label{proposition:appA-紧生成}
设 $X$ 为局部紧 Hausdorff 空间。则下述陈述等价：(i) $X$ 具有可数基；(ii) $X$ 可度量化且为 $\sigma$-紧。
\end{proposition}

\begin{definition}[Polish 空间]\label{definition:appA-Polish}
具有可数基的完全可度量化拓扑空间称为 Polish 空间。可分 Banach 空间以及具有可数基的局部紧 Hausdorff 空间都是 Polish 空间的例子。
\end{definition}

\section{连续函数空间}\label{section:appA-连续函数空间}

\begin{definition}[连续函数空间 $C$ 与 $C_b$]\label{definition:appA-C与Cb}
设 $X$ 为拓扑空间，$Y$ 为距离空间。记 $C(X;Y)$ 为连续映射 $f:X\to Y$ 构成的空间，$C_b(X;Y)\subseteq C(X;Y)$ 为有界连续函数构成的子空间。简记 $C(X)\equiv C(X;\mathbb{R})$，$C_b(X)\equiv C_b(X;\mathbb{R})$。若 $Y$ 为 Banach 空间，则 $C_b(X;Y)$ 赋以范数
\begin{align*}
f\mapsto\|f\|_\infty:=\sup_{x\in X}\|f(x)\|_Y
\end{align*}
后是 Banach 空间。
\end{definition}

回顾如下结果（例如见文献~\cite{ref107} 定理 I.5.1）。

\begin{proposition}\label{proposition:appA-Cb可分}\label{proposition:chA-C0可分}
\label{subsection:appA-可分}
\label{paragraph:appA-可分}
设 $X$ 为紧距离空间，$Y$ 为完备可分距离空间。则 $C_b(X;Y)$ 赋以范数 $\|\cdot\|_\infty$ 后是可分的。
\end{proposition}

\begin{definition}[上、下半连续]\label{definition:appA-半连续}
设 $X$ 为拓扑空间。函数 $f:X\to\mathbb{R}$ 称为上半连续（相应地，下半连续）的，若对每个 $\alpha\in\mathbb{R}$，集合 $\{f<\alpha\}$（相应地，$\{f>\alpha\}$）是开集。
\end{definition}

\begin{proposition}\label{proposition:appA-上半连续最大值}
设 $X$ 为可数紧拓扑空间，$f:X\to\mathbb{R}$ 为上半连续函数。则 $f$ 有上界且达到最大值。
\end{proposition}

\begin{definition}[一致连续]\label{definition:appA-一致连续}
设 $(X,d_X)$ 与 $(Y,d_Y)$ 为距离空间。映射 $f:X\to Y$ 称为一致连续的，若对每个 $\epsilon>0$，存在 $\delta>0$，使得对所有满足 $d_X(x,x')<\delta$ 的 $x,x'\in X$，都有 $d_Y(f(x),f(x'))<\epsilon$。若 $X$ 为紧空间，则每个 $f\in C(X;Y)$ 都有界且一致连续。
\end{definition}

\begin{definition}[紧支连续函数与 $C_0$]\label{definition:appA-Cc与C0}
记 $C_c(X)\subseteq C(X)$ 为紧支连续实值函数构成的子空间（回忆 $\operatorname{supp}f:=\{f\neq 0\}$）。记 $C_0(X)\subseteq C(X)$ 为具有如下性质的连续函数构成的子空间：对每个 $\epsilon>0$，存在紧集 $K\subseteq X$，使 $|f|$ 在 $K^c$ 上的限制小于 $\epsilon$。显然 $C_c(X)\subseteq C_0(X)$。$C_0(X)$ 赋以范数
\begin{align*}
f\mapsto\|f\|_\infty:=\sup_{x\in X}|f(x)|
\end{align*}
后是 Banach 空间。若 $X$ 为具有可数基的局部紧 Hausdorff 空间，则 $C_0(X)$ 赋以范数 $\|\cdot\|_\infty$ 后是可分 Banach 空间。
\end{definition}

\begin{proposition}[Dini 定理]\label{proposition:appA-Dini}
设 $X$ 为可数紧拓扑空间，$\{f_n\}$ 为一列非增的上半连续函数，且对所有 $x\in X$ 有 $\lim_{n\to\infty}f_n(x)=0$。则
\begin{align*}
\lim_{n\to\infty}\|f_n\|_\infty=0.
\end{align*}
\end{proposition}

\section{完全正则、正规与扩张定理}\label{section:appA-正则正规}

\begin{definition}[完全正则与正规]\label{definition:appA-完全正则正规}
拓扑空间 $X$ 称为完全正则的，若对所有 $x\in X$ 以及所有满足 $x\notin C$ 的闭集 $C\subseteq X$，存在 $f\in C(X)$ 使 $f(X)\subseteq[0,1]$，$f(x)=0$，且 $f|_C=1$。空间 $X$ 称为正规的，若对任意满足 $E\cap F=\emptyset$ 的闭集 $E,F\subset X$，存在 $E$ 的邻域 $U$ 与 $F$ 的邻域 $V$，使 $U\cap V=\emptyset$。每个距离空间都是完全正则且正规的。
\end{definition}

\begin{proposition}[Urysohn 引理]\label{proposition:appA-Urysohn}
拓扑空间 $X$ 是正规的，当且仅当对任意不相交的闭集 $E,F\subset X$，存在 $f\in C(X)$ 使 $f(X)\subseteq[0,1]$，$f|_E=0$，$f|_F=1$。
\end{proposition}

\begin{corollary}\label{corollary:appA-局部紧完全正则}
每个局部紧 Hausdorff 空间都是完全正则的。
\end{corollary}

\begin{proposition}[Tietze 扩张定理]\label{proposition:appA-Tietze}
设 $X$ 为正规拓扑空间，$F$ 为 $X$ 的闭子集，$f\in C(F)$。则存在 $g\in C(X)$ 使 $g|_F=f$。
\end{proposition}

如下引理表明：当 $X$ 局部紧时，$C_c(X)$ 在范数 $\|\cdot\|_\infty$ 下在 $C_0(X)$ 中稠密。

\begin{lemma}\label{lemma:appA-Urysohn}
\label{lemma:appA-光滑截断}
设 $X$ 为局部紧 Hausdorff 空间，$K\in\mathcal{K}(X)$，$A\in\mathscr{T}$ 且 $A\supset K$。则存在 $f\in C_c(X)$ 使 $f(X)\subseteq[0,1]$，$f|_K=1$，且 $\operatorname{supp}f\subseteq A$。
\end{lemma}

\begin{proof}
由于 $X$ 局部紧，存在 $B\in\mathscr{T}$ 使 $A\supseteq B\supseteq K$。由\cref{corollary:appA-局部紧完全正则}，$X$ 完全正则，因而对每个 $x\in K$，存在 $f_x\in C_c(X)$ 使 $f_x(X)\subseteq[0,1]$，$f_x(x)=1$，且 $f_x|_{B^c}=0$。由于 $K$ 紧，且每个集合 $\{f_x>1/2\}$（$x\in K$）是开集，故存在 $x_1,\dots,x_p$ 使
\begin{align*}
K\subseteq\bigcup_{k=1}^p\{f_{x_k}>1/2\}.
\end{align*}
对每个 $x\in X$，令
\begin{align*}
g(x):=2\max\{f_{x_1}(x),\dots,f_{x_p}(x)\},
\end{align*}
则 $g\in C_c(X)$，$g|_K>1$，且 $g|_{B^c}=0$。于是 $f:=\min\{1,g\}$ 即具有所述性质。
\end{proof}

\section{对偶空间上的弱$^*$拓扑}\label{section:appA-对偶弱星}

\begin{lemma}\label{lemma:appA-Ystar范数}
设 $Y$ 为可分赋范向量空间，$D\equiv\{y_k\mid k\in\mathbb{N}\}\subseteq Y$ 为可数稠密集。则由
\begin{align}
Y^*\ni y^*\mapsto |||y^*|||:=\sum_{k=1}^\infty\frac{|\langle y^*,y_k\rangle_{Y^*,Y}|}{2^k(1+\|y_k\|_Y)} \label{align:appA-Ystar范数}
\end{align}
定义的从 $Y^*$ 到 $[0,\infty)$ 的映射是 $Y^*$ 上的一个范数。
\end{lemma}

\begin{proposition}\label{proposition:appA-弱星拓扑}\label{proposition:appA-弱星可度量化}
设 $Y$ 为可分赋范向量空间，令
\begin{align}
V:=\{y^*\in Y^*\mid\|y^*\|_{Y^*}\leqslant 1\}. \label{align:appA-V单位球}
\end{align}
记 $\mathscr{T}_{w^*}$ 为 $Y^*$ 上的弱$^*$拓扑，$\widetilde{\mathscr{T}}$ 为 $Y^*$ 上由\eqref{align:appA-Ystar范数}中的范数 $|||\cdot|||$ 诱导的距离拓扑。则相对拓扑 $\mathscr{T}_{w^*}|_V$ 与 $\widetilde{\mathscr{T}}|_V$ 一致。

因而，对每个 $y\in Y$，
\begin{align*}
\lim_{n\to\infty}\langle y_n^*,y\rangle_{Y^*,Y}=\langle y^*,y\rangle_{Y^*,Y}
\quad\Longleftrightarrow\quad
\lim_{n\to\infty}|||y_n^*-y^*|||=0.
\end{align*}
\end{proposition}

\end{document}
